\section{苦涩的教训与3D的未来}

\subsection{Bitter Lesson 视角下的3D生成}

Henzler 引用了 Richard Sutton 的「苦涩的教训」（The Bitter Lesson）来审视3D生成领域的发展趋势：

\begin{quote}
\textit{假设数据和计算资源充裕，减少归纳偏置。}
\end{quote}

在这一框架下：
\begin{itemize}
    \item \textbf{真正的生成方法}（2D图像/视频模型）\textbf{能够扩展}（scale）：数据充裕、归纳偏置有限
    \item \textbf{3D模型无法同样扩展}：高质量3D数据极度稀缺
    \item 但3D模型在\textbf{成本和时间效率}上具有显著优势
\end{itemize}

\begin{importantbox}{3D在内容生成中的持续重要性}
虽然2D方法已能生成高度逼真的视频并展现了交互能力（如相机控制），但：

\textbf{1. 控制精度不足}：我们尚未获得对2D视频模型的完全、精确的控制能力。

\textbf{2. 部署困难}：这些大型模型目前无法在消费级设备上实时运行。

\textbf{3. 3D蒸馏不可或缺}：为了在移动端部署，仍然需要将生成结果蒸馏为高效的3D表示。

这些限制意味着3D建模在可预见的未来仍是内容生成流程中的关键一环。
\end{importantbox}

\subsection{World Labs 的启示}

World Labs 正在开发能够在\textbf{单个 GPU} 上实现\textbf{实时下一帧预测}的交互式3D模型。虽然结果尚不完美，但这是一个非常有前景的方向——预示着2D生成与3D表示之间的边界可能进一步模糊。

\subsection{本章小结}

3D生成的未来可能在于「2D负责质量、3D负责效率」的分工协作模式。3D方法的不可替代性在于其极低的渲染成本和在移动端的可部署性。

